% =====================================================================================================
% appB_wilson.tex —— 附录 B  Wilson 线：Feynman 规则、程函恒等式与尖点（writer W05）
% 源：FINAL_OUTLINE_v2 App B；research/c_wilson（DERIVATION.md §2--§3）；VF-WL-1/3/4/13，VF-CONV-23，VF-SCET-16；
%     验证脚本 writing/work/W05/s4_wilson.py、s5_cusp.py（日志 logs/）。
% =====================================================================================================
\chapter{Wilson 线：Feynman 规则、程函恒等式与尖点}\label{app:appB}

本附录证明第~3~章的\cref{prop:prelim-wilson-covariance}、\cref{prop:prelim-wilson-rules}（由\cref{prop:appB-rules}给出）与\cref{lem:prelim-eikonal-identity}（由\cref{lem:appB-shuffle}给出），
给出半无限 Wilson 线的 Ward 恒等式（供附录~F 使用）、非阿贝尔指数化的全阶论证（\cref{prop:appB-webs}）与第~5~章所用的单圈 cusp 反常维度（\cref{prop:appB-cusp-oneloop}）。

% -----------------------------------------------------------------------------------------------------
\section{定义与约定}\label{sec:appB-definition}

沿路径 $z(s)$（$s\in[0,1]$，$z(0)=a$，$z(1)=b$）的 Wilson 线（\cref{def:prelim-wilson}）为
\begin{equation}\label{eq:appB-W}
W[b,a]=P\exp\Big(\ii g\int_a^b\dd z\cdot A\Big)=\sum_{m=0}^\infty(\ii g)^m\!\!\int\limits_{0<s_1<\dots<s_m<1}\!\!\dd s_1\cdots\dd s_m\;
\dot z\cdot A(z(s_m))\cdots\dot z\cdot A(z(s_1)),
\end{equation}
$A\equiv A^at^a$，离 $b$ 近的点在左；$W(s)\equiv W[z(s),a]$ 是 $\frac{\dd}{\dd s}W=\ii g\,\dot z\cdot A(z(s))\,W$、$W(0)=1$ 的唯一解。

\emph{\cref{prop:prelim-wilson-covariance}的证明。}由于 $W[c,b]W[b,a]$ 与 $W[c,a]$ 满足同一方程与同一初值，二者相等；取 $c=a$ 得 $W[a,b]W[b,a]=1$，对厄米的 $A$ 有 $W[b,a]^\dagger=W[a,b]$。
在 $D=\partial-\ii gA$ 的规范变换 $A'=UAU^{-1}-\frac{\ii}g(\partial U)U^{-1}$ 下，$V(s)=U(z(s))W(s)U^{-1}(a)$ 满足
\[
\frac{\dd V}{\dd s}=\Big[\frac{\dd U}{\dd s}U^{-1}+\ii g\,U\dot z\cdot AU^{-1}\Big]V=\ii g\,\dot z\cdot A'\,V,\qquad V(0)=1,
\]
故 $W[b,a]\to U(b)W[b,a]U^{-1}(a)$（非齐次项反号则不成立；SU(3) 中数值核对）。% check: s4_wilson.py (d)
伴随线 $W^{(A)}_{ab}=2\Tr[t^aW[b,a]t^bW[a,b]]$ 即 $W t^bW^{-1}=t^aW^{(A)}_{ab}$。

\emph{约定。}(C1) $A(x)=\int\frac{\dd^4k}{(2\pi)^4}\ee^{-\ii k\cdot x}\tilde A(k)$，$k$ 流入 Wilson 线；(C2) $n^0>0$，$W[\infty n,y]=\lim_{L\to\infty}W[y+Ln,y]$ 等，端点 $y\ne0$ 时乘 $\ee^{-\ii y\cdot K_{\rm tot}}$；
(C3) 绝热正规子：半无限线上的场乘以 $\ee^{-\epsilon s}$（$s$ 为到有限端点的仿射距离），最后 $\epsilon\to0^+$。Collins 用 $D=\partial+\ii g_0tA$\cite[式~(2.2)]{Collins2011}，故 $g_{\rm Collins}=-g$，
其 Wilson 线 $P\exp(-\ii g_0\int n\cdot A)$\cite[式~(13.40)]{Collins2011}作为算符与我们的相同；他的图~7.10 与式~(11.4)--(11.6) 须按作者勘误使用（附录~Q 第~19~条）。

% -----------------------------------------------------------------------------------------------------
\section{动量空间 Feynman 规则}\label{sec:appB-rules}

本节给出第~3~章\cref{prop:prelim-wilson-rules}的证明。

\begin{proposition}[Wilson 线的 Feynman 规则 \tagR]\label{prop:appB-rules}
（$R_{\mathrm{def}}$：由算符定义直接推出。）按 (C1)--(C3)，直半无限 Wilson 线中连接 $m$ 个胶子（动量 $k_j$ 流入、指标 $\mu_j$、色 $a_j$；记 $n\cdot\tilde A_j\equiv n^{\mu_j}\tilde A^{a_j}_{\mu_j}(k_j)t^{a_j}$，按所写次序相乘）的项为
\begin{align}
W^{(m)}[0,-\infty n]&=\int\!\prod_j\frac{\dd^4k_j}{(2\pi)^4}\,(\ii g)^m\,n\cdot\tilde A_m\cdots n\cdot\tilde A_1\prod_{j=1}^m\frac{\ii}{n\cdot(k_1+\dots+k_j)+\ii j\epsilon},\label{eq:appB-past}\\
W^{(m)}[\infty n,0]&=\int\!\prod_j\frac{\dd^4k_j}{(2\pi)^4}\,(\ii g)^m\,n\cdot\tilde A_m\cdots n\cdot\tilde A_1\prod_{j=1}^m\frac{-\ii}{n\cdot(k_j+\dots+k_m)-\ii(m-j+1)\epsilon},\label{eq:appB-future}\\
W^{(m)}[-\infty n,0]&=\int\!\prod_j\frac{\dd^4k_j}{(2\pi)^4}\,(-\ii g)^m\,n\cdot\tilde A_m\cdots n\cdot\tilde A_1\prod_{j=1}^m\frac{\ii}{n\cdot(k_j+\dots+k_m)+\ii(m-j+1)\epsilon},\label{eq:appB-pastc}\\
W^{(m)}[0,\infty n]&=\int\!\prod_j\frac{\dd^4k_j}{(2\pi)^4}\,(-\ii g)^m\,n\cdot\tilde A_m\cdots n\cdot\tilde A_1\prod_{j=1}^m\frac{-\ii}{n\cdot(k_1+\dots+k_j)-\ii j\epsilon},\label{eq:appB-futc}
\end{align}
% VF-WL-1
其中胶子按路径次序编号：式~\eqref{eq:appB-past} 中 $k_1$ 最靠近 $-\infty$，式~\eqref{eq:appB-future} 与 \eqref{eq:appB-pastc} 中 $k_1$ 最靠近 $0$，式~\eqref{eq:appB-futc} 中 $k_1$ 最靠近 $+\infty$。
等价地（普适规则）：每个连接点一个顶点 $\ii g\,t^av^\mu$（$v$ 为沿取向的切矢量），除伸向无穷远的一段外每段一个 $\ii/(v\cdot\ell+\ii0)$（$\ell$ 沿取向流动，伸向无穷远的一段 $\ell=0$），
色矩阵按路径次序，有限 $\epsilon$ 下某段的 $\ii0$ 为 $\ii\epsilon$ 乘以该段与无穷远端之间的连接点数。单胶子项与“未来减过去”为
\begin{equation}\label{eq:appB-onegluon}
W[\infty n,0]\ni\frac{g\,t^an\cdot\tilde A^a(k)}{n\cdot k-\ii0},\qquad W[-\infty n,0]\ni\frac{g\,t^an\cdot\tilde A^a(k)}{n\cdot k+\ii0},\qquad
\text{差}=2\pi\ii\,g\,t^an\cdot\tilde A^a(k)\,\delta(n\cdot k).
\end{equation}
\uses{def:prelim-wilson}
\end{proposition}

\begin{proof}
以 $W[0,-\infty n]=P\exp(\ii g\int_{-\infty}^0\dd\lambda\,n\cdot A(\lambda n))$ 为例：代入 (C1) 并乘绝热因子，$m$ 阶项的参数积分为
$\int_{-\infty<\lambda_1<\dots<\lambda_m<0}\prod_j\dd\lambda_j\,\ee^{-\ii\lambda_j(n\cdot k_j+\ii\epsilon)}$。先积 $\lambda_1$ 得 $\frac{\ii}{n\cdot k_1+\ii\epsilon}\ee^{-\ii\lambda_2(n\cdot k_1+\ii\epsilon)}$，
它与下一个相因子合并为 $\ee^{-\ii\lambda_2(n\cdot(k_1+k_2)+2\ii\epsilon)}$；归纳地第 $j$ 次积分给出 $\ii/[n\cdot(k_1+\dots+k_j)+\ii j\epsilon]$，$\lambda_m$ 积到 $0$。色矩阵的次序即式~\eqref{eq:appB-W} 的路径次序。
$W[\infty n,0]$ 的区域为 $0<\lambda_1<\dots<\lambda_m<\infty$、阻尼 $\ee^{-\epsilon\lambda_j}$，先积 $\lambda_m$，得式~\eqref{eq:appB-future}；$W[-\infty n,0]$、$W[0,\infty n]$ 的取向为 $-n$，
每个顶点多一个负号，动量仍从无穷远端累加。普适规则是这四式的改写（例如 $W[\infty n,0]$ 第 $j$ 个胶子前一段沿取向的动量为 $-(k_j+\dots+k_m)$）。
式~\eqref{eq:appB-onegluon} 是 $m=1$ 的情形，差由 $\frac1{x-\ii0}-\frac1{x+\ii0}=2\pi\ii\delta(x)$ 给出（即 $\ii g\int_{-\infty}^\infty\dd\lambda\,n\cdot A(\lambda n)$ 的 Fourier 变换）。
四种线的 $m\le3$ 项经符号计算验证；把式~\eqref{eq:appB-W} 在 SU(3) 背景场中数值求解，$g^1$ 至 $g^4$ 的系数与规则相符（残差按 $g^5$ 标度）。% check: s4_wilson.py (a)(b)
\end{proof}
% -----------------------------------------------------------------------------------------------------
\section{程函恒等式与 Wilson 线的 Ward 恒等式}\label{sec:appB-eikonal}

本节给出第~3~章\cref{lem:prelim-eikonal-identity}的证明。

\begin{lemma}[程函恒等式（洗牌积）]\label{lem:appB-shuffle}
(i) 设复数 $y_1,\dots,y_m$ 的任意非空子集之和都不为零，则
\begin{equation}\label{eq:appB-shuffle}
\sum_{\sigma\in S_m}\prod_{j=1}^m\frac{1}{y_{\sigma(1)}+\dots+y_{\sigma(j)}}=\prod_{j=1}^m\frac1{y_j}.
% VF-WL-4
\end{equation}
更一般地，对两个保持内部次序的序列 $(q_1,\dots,q_p)$ 与 $(r_1,\dots,r_q)$，记 $E(x_1,\dots,x_l)=\prod_{j=1}^l(x_1+\dots+x_j)^{-1}$，则对全部 $\binom{p+q}{p}$ 个洗牌 $\sigma$ 求和有
$\sum_\sigma E(\sigma)=E(q_1,\dots,q_p)\,E(r_1,\dots,r_q)$。
(ii) 取 $y_j=n\cdot k_j+\ii\epsilon$（$\epsilon>0$），则左边第 $j$ 个分母为 $n\cdot(k_{\sigma(1)}+\dots+k_{\sigma(j)})+\ii j\epsilon$，正是绝热正规子 (C3) 的累积形式：
程函恒等式 $\sum_\sigma\prod_j\frac{\ii}{n\cdot(k_{\sigma(1)}+\dots+k_{\sigma(j)})+\ii j\epsilon}=\prod_j\frac{\ii}{n\cdot k_j+\ii\epsilon}$ 在\emph{有限} $\epsilon$ 下严格成立（指向未来的线同理）。
(iii) 若每个程函分母都加\emph{同一个} $+\ii\epsilon$，恒等式只在 $\epsilon\to0$ 时按分布意义成立：$m=2$ 时（$a_j=n\cdot k_j$）
\[
\prod_j\frac{1}{a_j+\ii\epsilon}-\sum_\sigma\prod_j\frac{1}{a_{\sigma(1)}+\dots+a_{\sigma(j)}+\ii\epsilon}=\frac{-\ii\epsilon}{(a_1+\ii\epsilon)(a_2+\ii\epsilon)(a_1+a_2+\ii\epsilon)},
\]
它与 $\frac1{(1+a_1^2)(1+a_2^2)}$ 的配对为 $\frac{\pi^2\epsilon}{(1+\epsilon)^2(2+\epsilon)}$，只在 $\epsilon\to0$ 时为零。
% uses: none
\end{lemma}

\begin{proof}
(i) 对 $m$ 归纳。按最后一个元素 $\sigma(m)=l$ 分组：末因子 $1/(y_1+\dots+y_m)$ 公共，其余由归纳假设等于 $\prod_{j\ne l}1/y_j$，故左边
$=\frac{1}{\sum_jy_j}\sum_l\prod_{j\ne l}\frac1{y_j}=\prod_j\frac1{y_j}$。洗牌积同样对 $p+q$ 归纳（最后一个元素是 $q_p$ 或 $r_q$）。等价地，$\prod_j\frac{\ii}{y_j}=\prod_j\int_{-\infty}^0\dd\lambda_j\ee^{-\ii\lambda_jy_j}$
（$\Imag y_j>0$），而 $(-\infty,0]^m$ 是 $m!$ 个单形的不交并，每个单形给出左边一项：阿贝尔 Wilson 线是普通指数，$W^{(m)}=(W^{(1)})^m/m!$。
(ii) 由 (i)：$y_{\sigma(1)}+\dots+y_{\sigma(j)}$ 的虚部为 $j\epsilon$。(iii) 通分即得；配对用留数：对 $a_1$ 闭合于上半平面只有 $a_1=\ii$，给出 $\frac{-\ii\pi}{(1+\epsilon)(a_2+\ii(1+\epsilon))}$，
再对 $a_2$ 同样得 $\frac{-\pi}{(2+\epsilon)(1+\epsilon)}$，乘以 $-\ii\epsilon$ 即得（亦经数值核对）。% check: writing/work/W05/s4_wilson.py (e)
\end{proof}

即：只要每个累积分母的正规子等于已吸收胶子的正规子之和，恒等式在有限正规子下就严格成立——场的绝热阻尼正给出这种累积形式；朴素的统一 $+\ii\epsilon$（$\delta$ 正规化）不满足这一条件，恒等式随之失效（(iii)；$m=2$ 时这一条件也是必要的）
（参见\cite[式~(5)--(6)]{Echevarria:2015byo}；阿贝尔恒等式本身是经典的）。

\begin{remark}[半无限 Wilson 线的 Ward 恒等式]\label{rem:appB-ward}
在 $W[0,-\infty n]$ 的 $m$ 胶子项中（色矩阵 $T_1,\dots,T_m$ 已含顶点因子，$T_1$ 最靠近 $-\infty$；记该项为 $Y$），把一个动量为 $q$、极化为 $q^\mu$、色为 $T_q$ 的标量极化胶子插入全部 $m+1$ 个位置并求和，得
\begin{equation}\label{eq:appB-ward}
\sum_{p=0}^{m}(\text{插在第 }p\text{ 个胶子之后})=-g\,T_q\,Y+g\sum_{p=0}^{m-1}Z_p\;T_m\cdots T_{p+2}\,[T_q,T_{p+1}]\,T_p\cdots T_1 ,
% VF-WL-13
\end{equation}
其中 $Z_p$ 是第 $p+1$ 个胶子携带 $k_{p+1}+q$ 时的程函传播子之积。证明：内部插入用程函 Ward 恒等式
$\frac{\ii}{n\cdot(\ell+q)}(\ii g\,n\cdot q)\frac{\ii}{n\cdot\ell}=g[\frac{\ii}{n\cdot(\ell+q)}-\frac{\ii}{n\cdot\ell}]$ 拆开，最靠近 $-\infty$ 的插入给出 $-g$，相邻位置的项逐对相消，只剩相邻色序之差与端点项；
阿贝尔情形和为 $-gY$。对易子项与该胶子经三胶子顶点的图相消（附录~F）。若标量极化也带同样的绝热移动，式~\eqref{eq:appB-ward} 在有限 $\epsilon$ 下严格成立（以 SU(3) 生成元数值核对至 $m=4$）。% check: s4_wilson.py (f)
\uses{prop:appB-rules}
\end{remark}

\begin{proposition}[非阿贝尔指数化 \tagA]\label{prop:appB-webs}
在一点相遇的 Wilson 线的真空关联函数（及其在背景场中的期望值）可以写成 web 之和的指数：$\langle W\rangle=\exp\big[\sum_w\tilde C(w)\,\mathcal F(w)\big]$，其中 $\mathcal F(w)$ 是 web 的运动学积分，
$\tilde C(w)$ 是修正的色因子；对两条线，web 是交换胶子在两条线上的连接不能被分成两个独立子集的图，其修正色因子是“最大非阿贝尔”部分\cite{Gatheral1983,FrenkelTaylor1984}。
\uses{prop:appB-rules,lem:appB-shuffle}
\end{proposition}

\begin{proofsketch}
阿贝尔情形：由\cref{prop:appB-rules}与洗牌积（\cref{lem:appB-shuffle}），独立的连接组之间因子化，$\langle W\rangle$ 是连通光子子图之和的指数。非阿贝尔情形：把每个图的色因子按交换次序写成对称部分与对易子之和，
对称部分用洗牌积重组进指数，对易子部分留下修正色因子 $\tilde C(w)$（两条线时为“最大非阿贝尔”部分）；$O(g^2)$ 时即\cref{prop:sudakov-webs}(b) 的严格恒等式（指针），两圈交叉梯形图的修正色因子为 $-\frac12C_FC_A$。
完整论证见\cite{Gatheral1983,FrenkelTaylor1984}。
\end{proofsketch}

\paragraph{地位与缺口}全阶的组合学论证（[A]），对两条线在文献中给出\cite{Gatheral1983,FrenkelTaylor1984}；在 $O(g^2)$ 是严格的代数恒等式（\cref{prop:sudakov-webs}(b)，[R, $R_{\mathrm{def}}$]）。
多条线的 web 及其重整化本书不需要。所需的有限正规子必须是绝热型的（\cref{lem:appB-shuffle}(iii)）。
% -----------------------------------------------------------------------------------------------------
\section{单圈尖点反常维度}\label{sec:appB-cusp}

\begin{proposition}[单圈尖点反常维度 \tagR]\label{prop:appB-cusp-oneloop}
（$R_{\mathrm{PT}}$：在重整化微扰论中对红外正规化的部分子态严格。）在 $O(\alpha_s)$，对散射组态的折线
$W_\beta=\frac1{N_c}\langle0|\Tr\,T\{W[\infty v_1,0]\,W[0,-\infty v_2]\}|0\rangle$（$v_i^2=1$，$v_i^0>0$，$\cosh\beta=v_1\cdot v_2$；Feynman 规范，$D=4-2\varepsilon$，红外由绝热阻尼 $\ee^{-\delta|\lambda|}$ 正规化），
\begin{equation}\label{eq:appB-cusp-pole}
W_\beta=1-\frac{\alpha_sC_F}{2\pi\varepsilon}\,(\beta\coth\beta-1)+(\varepsilon\text{ 的正则项})+O(\alpha_s^2),\qquad
\Gamma_{\rm cusp}(\beta)=\frac{\alpha_sC_F}{\pi}\,(\beta\coth\beta-1)+O(\alpha_s^2),
% VF-CONV-23
\end{equation}
其中 $\Gamma_{\rm cusp}$ 按\cref{def:sudakov-cusp}定义（$\MSbar$）。$\beta\coth\beta$ 来自两线之间的交换，$-1$ 来自两条线的自能；$\Gamma(0)=0$；
$\beta\to\infty$ 时 $\Gamma=\frac{\alpha_sC_F}\pi(\beta-1)+O(\beta\ee^{-2\beta})$，线性项的系数 $\frac{\alpha_sC_F}\pi=\frac{\alpha_s}{4\pi}\Gamma_0$ 即单圈的类光 cusp 反常维度，
而 $v_i=(\cosh y_i,0,0,\sinh y_i)$ 时 $\beta=|y_1-y_2|$。
\uses{prop:appB-rules,def:prelim-dr-conventions,def:prelim-msbar,def:prelim-feynman-rules,def:sudakov-cusp}
\end{proposition}

\begin{proof}
\emph{第一步：图。}$O(g^2)$ 有三个图：两条线之间的单胶子交换，与每条线自身的 $O(g^2)$ 项（自能）。由\cref{prop:appB-rules}的单胶子项与 Schwinger 表示
$\frac{-\ii}{v\cdot k-\ii\delta}=\int_0^\infty\dd l\,\ee^{-\ii l(v\cdot k-\ii\delta)}$、$\frac{\ii}{v\cdot k+\ii\delta}=\int_{-\infty}^0\dd\lambda\,\ee^{-\ii\lambda(v\cdot k+\ii\delta)}$，
$W^{(1)}[\infty v_1,0]=\ii g\int_0^\infty\dd l\,\ee^{-\delta l}\,v_1\cdot A(lv_1)$，$W^{(1)}[0,-\infty v_2]=\ii g\int_{-\infty}^0\dd\lambda\,\ee^{\delta\lambda}\,v_2\cdot A(\lambda v_2)$。
Feynman 规范中 $\langle0|TA^a_\mu(x)A^b_\nu(y)|0\rangle=-g_{\mu\nu}\delta^{ab}D_F(x-y)$（\cref{def:prelim-feynman-rules}），
$D_F(x)=\int\frac{\dd^Dk}{(2\pi)^D}\frac{\ii\,\ee^{-\ii k\cdot x}}{k^2+\ii0}=\frac{\Gamma(1-\varepsilon)}{4\pi^{2-\varepsilon}}(-x^2+\ii0)^{-1+\varepsilon}$（Wick 转动后的欧氏 Fourier 变换，\cref{def:prelim-dr-conventions}）。
色因子 $\frac1{N_c}\Tr(t^at^a)=C_F$。令 $s=-\lambda$，
\[
\begin{aligned}
E&=g^2C_F\cosh\beta\int_0^\infty\!\dd l\int_0^\infty\!\dd s\,\ee^{-\delta(l+s)}D_F(lv_1+sv_2),\\
\Sigma_i&=g^2C_F\!\!\int\limits_{0<l_1<l_2}\!\!\dd l_1\dd l_2\,\ee^{-\delta(l_1+l_2)}D_F\big((l_2-l_1)v_i\big)=\frac{g^2C_F}{2\delta}\int_0^\infty\!\dd\tau\,\ee^{-\delta\tau}D_F(\tau v_i)
\end{aligned}
\]
（$E$ 的符号：$(\ii g)^2\cdot(-v_1\cdot v_2)=g^2\cosh\beta$；$\Sigma_i$ 中 $(\ii g)^2\cdot(-v_i^2)=g^2$，令 $\tau=l_2-l_1$ 后对 $l_1$ 积分得 $1/(2\delta)$）。

\emph{第二步：交换图。}令 $l=\rho u$、$s=\rho(1-u)$（$\dd l\,\dd s=\rho\,\dd\rho\,\dd u$），$(lv_1+sv_2)^2=\rho^2N(u)$，$N(u)=u^2+(1-u)^2+2u(1-u)\cosh\beta=1+2u(1-u)(\cosh\beta-1)\ge1$。
间隔是类时的，$(-\rho^2N+\ii0)^{-1+\varepsilon}=-\ee^{\ii\pi\varepsilon}(\rho^2N)^{-1+\varepsilon}$，于是
\[
\begin{aligned}
E&=-g^2C_F\cosh\beta\,\frac{\ee^{\ii\pi\varepsilon}\Gamma(1-\varepsilon)}{4\pi^{2-\varepsilon}}\int_0^\infty\!\dd\rho\,\rho^{-1+2\varepsilon}\ee^{-\delta\rho}\int_0^1\!\dd u\,N(u)^{-1+\varepsilon}\\
&=-g^2C_F\cosh\beta\,\frac{\ee^{\ii\pi\varepsilon}\Gamma(1-\varepsilon)\Gamma(2\varepsilon)}{4\pi^{2-\varepsilon}\,\delta^{2\varepsilon}}\int_0^1\!\dd u\,N(u)^{-1+\varepsilon}.
\end{aligned}
\]
$\rho\to0$ 处的发散（尖点处的紫外发散）给出极点 $\Gamma(2\varepsilon)=\frac1{2\varepsilon}+O(1)$。令 $w=2u-1$，$N=\frac12(a^2-b^2w^2)$，$a^2=\cosh\beta+1$，$b^2=\cosh\beta-1$，$ab=\sinh\beta$，
\[
\int_0^1\frac{\dd u}{N(u)}=\int_{-1}^{1}\frac{\dd w}{a^2-b^2w^2}=\frac{1}{ab}\ln\frac{a+b}{a-b}=\frac{\beta}{\sinh\beta}\qquad\Big(\frac{a+b}{a-b}=\frac{\cosh\frac\beta2+\sinh\frac\beta2}{\cosh\frac\beta2-\sinh\frac\beta2}=\ee^\beta\Big).
\]
故 $E$ 的极点为 $-\frac{g^2C_F}{8\pi^2\varepsilon}\beta\coth\beta=-\frac{\alpha_sC_F}{2\pi\varepsilon}\beta\coth\beta$。

\emph{第三步：自能。}$D_F(\tau v)=-\ee^{\ii\pi\varepsilon}\frac{\Gamma(1-\varepsilon)}{4\pi^{2-\varepsilon}}\tau^{-2+2\varepsilon}$，而
$\int_0^\infty\dd\tau\,\tau^{-2+2\varepsilon}\ee^{-\delta\tau}=\Gamma(-1+2\varepsilon)\,\delta^{1-2\varepsilon}$
（在 $\Real\varepsilon>\frac12$ 收敛，再解析延拓：线性紫外发散在维数正规化中不留痕迹），故
\[
\Sigma_i=-g^2C_F\,\frac{\ee^{\ii\pi\varepsilon}\Gamma(1-\varepsilon)\Gamma(-1+2\varepsilon)}{8\pi^{2-\varepsilon}\,\delta^{2\varepsilon}} .
\]
由 $\Gamma(-1+2\varepsilon)=-\Gamma(2\varepsilon)/(1-2\varepsilon)=-\frac1{2\varepsilon}+O(1)$，
每条线的自能极点为 $+\frac{g^2C_F}{16\pi^2\varepsilon}=+\frac{\alpha_sC_F}{4\pi\varepsilon}$；它是半无限线在有限端点（尖点）处的对数紫外发散。

\emph{第四步：反常维度。}代入 $g_0^2=4\pi\alpha_s\mu^{2\varepsilon}\ee^{\varepsilon\gamma_E}(4\pi)^{-\varepsilon}$（\cref{def:prelim-msbar}），$\mu^{2\varepsilon}\delta^{-2\varepsilon}=(\mu^2/\delta^2)^\varepsilon$，
三图之和的极点为 $-\frac{\alpha_sC_F}{2\pi\varepsilon}(\beta\coth\beta-1)$，即式~\eqref{eq:appB-cusp-pole} 的第一式；故 $Z_{\rm cusp}=1-\frac{\alpha_sC_F}{2\pi\varepsilon}(\beta\coth\beta-1)$，
由 $\frac{\dd\alpha_s}{\dd\ln\mu}=-2\varepsilon\alpha_s+O(\alpha_s^2)$，$\Gamma_{\rm cusp}=\frac{\dd\ln Z_{\rm cusp}}{\dd\ln\mu}=\frac{\alpha_sC_F}\pi(\beta\coth\beta-1)$。
等价地，$W^{\rm ren}_\beta=1-\frac{\alpha_sC_F}{2\pi}(\beta\coth\beta-1)\ln\frac{\mu^2}{\delta^2}+\text{常数}$，$-\frac{\dd\ln W^{\rm ren}_\beta}{\dd\ln\mu}=\Gamma_{\rm cusp}$。

\emph{第五步：极限。}$\beta\coth\beta-1=\frac{\beta^2}3+O(\beta^4)$，$\Gamma(0)=0$：直线没有尖点发散，交换图的 $\int\dd u/N=1$ 恰好抵消两段在连接点的端点发散；
$\beta\to\infty$ 时 $\coth\beta=1+2\ee^{-2\beta}+\cdots$；$v_1\cdot v_2=\cosh(y_1-y_2)$。各积分与常数因子都经符号与数值计算核对。% check: s5_cusp.py
\end{proof}

\paragraph{两圈系数与文献注记。}在 $\Gamma_{\rm cusp}=\sum_n(\frac{\alpha_s}{4\pi})^{n+1}\Gamma_n$ 的归一化下，类光两圈系数为
$\Gamma_1=4C_F[(\frac{67}9-\frac{\pi^2}3)C_A-\frac{20}9T_Fn_f]$\cite{KorchemskyRadyushkin1987}，$n_f=5$ 时 $\Gamma_1=36.8436$；\cite[式~(A.10)]{BecherSchwartz2008}（JHEP 版与 arXiv v1、v2）在 $T_Fn_f$ 前多印了一个 $C_F$，按字面得到 $26.97$（附录~Q 第~1~条）。% VF-CONV-23, VF-SCET-16
式~\eqref{eq:appB-cusp-pole} 包含两条线的自能；Becher 与 Neubert 的有质量部分子形式\cite[式~(13)--(15)]{Becher:2009kw}只把 $\beta\coth\beta$ 记入尖点项、把 $-1$ 归入 $2\gamma^Q$（$\gamma^Q_0=-2C_F$）。
单圈结果见\cite{Polyakov1980,BrandtNeriSato1981}；产生组态（$\beta\to\beta-\ii\pi$）见第~5~章\cref{prop:sudakov-coulomb}。
